\documentclass[11pt]{article}
\usepackage{listings}
\usepackage{xcolor}
\usepackage{multirow}
\usepackage{multicol}
\usepackage{xparse}
\usepackage{amsmath}
\usepackage[BoldFont]{xeCJK}
\usepackage[most]{tcolorbox}
\usepackage[T1]{fontenc}

\usepackage{lmodern}

\graphicspath{ {./image/} }
\renewcommand\thesection{Question \ \arabic{section}}
\renewcommand\thesubsection{\arabic{section}-(\alph{subsection})}
\usepackage{hyperref}
\hypersetup{
    colorlinks=true,
    linkcolor=blue,
    filecolor=magenta,      
    urlcolor=cyan,
    pdftitle={習題6 },
    pdfpagemode=FullScreen,
}
\usepackage[a4paper, margin={3cm}]{geometry}
\linespread{1.15}
\lstset{language=C,keywordstyle={\bfseries \color{blue}}}
	

\title{習題6 \\\Large 總體經濟學原理與實習}
\author{B10703033\ 財金一\ 邱奕璉
\\ B10703103\ 財金一\ 許羽婷
\\ B10705007\ 資管一\ 劉冠甫
\\ B10705047\ 資管一\ 邱秉辰}
\date{April 22, 2022}

\setlength{\parindent}{0em}
\setlength{\parskip}{0.5em}

\setCJKmainfont{GenRyuMin TW}
\begin{document}



\maketitle

\section{}
\subsection{}
\begin{enumerate}
    \item The effects of supply shocks.
    \item The extent to which the economy is operating above or below capacity.
    \item People’s expectation of inflation.
\end{enumerate}
\subsection{}
The effect of Supply shocks.
\subsection{}
People’s expectations of inflation \\ \\
\textbf{Expectations are the trickiest piece of the inflation equation.} Measurement is one problem. A survey by the University of Michigan suggests that consumers expect average annual inflation of 3.1\% over the next five years; market-based measures imply a rate of about 2.6\% over the same period, before falling to about 2.2\% over the subsequent five years. That is above the Fed’s 2\% target but still well short of a 1970s-style rerun. Perhaps punters are less worried about price pressures. Or perhaps they trust in the Fed’s commitment to price stability.
\section{}
\subsection{}
本行重貼現率、擔保放款融通利率及短期融通利率分別維持年息1.125\%、1.5\%及3.375\%。
\subsection{}
金融業隔夜拆款利率:0.082\%
台灣重貼現利率 > 1\%
重貼現率 > 隔夜拆款利率 $\implies$ 金融機構無動機向央行借錢。
\subsection{}
0.08\% ; 美國利率較低，對資金較寬鬆。
\section{}
\subsection{}
\begin{figure}[h]
    \centering
    \includegraphics[scale = 0.15]{3a.png}
\end{figure}
台灣的利率都在南韓之下，代表台灣的貨幣政策較寬鬆。
\subsection{}
台灣的貨幣政策較寬鬆。
\section{}
\subsection{}
\begin{itemize}
    \item 台北市：16.29
    \item 新北市：12.52
\end{itemize}
民國101年第2季(含)以前採用財團法人聯合徵信中心擔保品建估總價，民國101年第3季起採用不動產成交案件實際資訊申報登錄、財政部財稅資料中心個人綜合所得稅申報資料，分別統計中位數房價、家戶年可支配所得中位數，再以中位數房價除以家戶年可支配所得中位數得出。
\subsection{}
台北市：
\begin{itemize}
    \item 2002 Q1: 6.06
    \item 2021 Q4: 16.29
\end{itemize}
上漲為 $16.29 / 6.06 \approx 2.688$ 倍
\subsection{}
$$
\frac{1000 + \frac{1000}{64}}{1000} - 1 = 1.56\%
$$
\subsection{}
$$
\frac{1000(0.02)}{1000} = 2\%
$$
\subsection{}
某甲會把錢存入定期存款，因為具有較高的報酬率。
\subsection{}
$$
\frac{1100 + \frac{1000}{64}}{1000} - 1 = 11.56\% 
$$

\end{document}


